	\documentclass[a4paper]{article}

\usepackage[spanish]{babel}
\usepackage{graphicx}
\usepackage{amsmath, amssymb}
\usepackage[margin=2cm]{geometry}
\usepackage{fancyhdr}
\usepackage{enumerate}
\usepackage[shortlabels]{enumitem}
\usepackage{parskip}
\usepackage[hidelinks]{hyperref}
\usepackage{float}
\usepackage{booktabs}



% cabecera
\pagestyle{fancy}
\fancyhead[l]{Mecánica Celeste}
\fancyhead[c]{Solución del Taller - Clase 2}
\fancyhead[r]{\today}
\fancyfoot[c]{\thepage}
\renewcommand{\headrulewidth}{0.2pt} % linea horizontal

\begin{document}

\begin{titlepage}
  \centering
  {\Large Universidad de Antioquia\par}
  \vspace{2cm}
  {\huge\bfseries Mecánica Celeste\par}
  \vspace{0.4cm}
  {\Large Clase 2 - Taller\par}
  \vspace{2.5cm}
  {\large Autor:\par}
  \vspace{0.2cm}
  {\large Daniel Soto Villada\par}
  {\large Estudiante de Astronomía\par}
  \vfill
  {\large \today\par}
\end{titlepage}

\newpage
\tableofcontents
\newpage

\section{Problema \# 1}

Sea el vector de posición definido en coordenadas cartesianas como $\vec{r} = (x, y, z)$ y su magnitud escalar definida como $r = |\vec{r}| = (x^2 + y^2 + z^2)^{1/2}$.

\begin{enumerate}[a)]
    \item Calcule las derivadas parciales $\frac{\partial r}{\partial x}$, $\frac{\partial r}{\partial y}$ y $\frac{\partial r}{\partial z}$.
    \item A partir de lo anterior, demuestre que el gradiente de la magnitud del vector de posición es $\nabla r = \frac{\vec{r}}{r}$ y verifique que la magnitud de este gradiente es $|\nabla r| = 1$ para todo $r \neq 0$.
\end{enumerate}

\section*{Solución Problema \# 1}

Este problema trata sobre el análisis básico del \textbf{vector de posición} en el espacio euclidiano tridimensional. El objetivo es comprender cómo varía la distancia radial $r$ desde el origen respecto a las coordenadas cartesianas y cómo estas variaciones se combinan para formar el \textbf{operador gradiente} aplicado a una función puramente radial.

\subsection{Parte (a): Cálculo de las derivadas parciales}

Para calcular las derivadas parciales de la función escalar $r(x, y, z) = (x^2 + y^2 + z^2)^{1/2}$, debemos aplicar la \textbf{regla de la cadena} para derivadas multivariadas.

\subsubsection{Derivada respecto a $x$}
Consideramos a las variables $y$ y $z$ como constantes durante la diferenciación:

$$ \frac{\partial r}{\partial x} = \frac{\partial}{\partial x} (x^2 + y^2 + z^2)^{1/2} $$

Aplicando la regla de la potencia y la derivada interna:

$$ \frac{\partial r}{\partial x} = \frac{1}{2} (x^2 + y^2 + z^2)^{-1/2} \cdot \frac{\partial}{\partial x} (x^2 + y^2 + z^2) $$

$$ \frac{\partial r}{\partial x} = \frac{1}{2} (x^2 + y^2 + z^2)^{-1/2} \cdot (2x) $$

Simplificando el factor $2$ y reescribiendo la potencia negativa en el denominador:

$$ \frac{\partial r}{\partial x} = \frac{x}{(x^2 + y^2 + z^2)^{1/2}} $$
$$\frac{\partial r}{\partial x} = \frac{x}{r} $$

\subsubsection{Derivadas respecto a $y$ y $z$}
Dada la \textbf{simetría esférica} de la función $r$ (es decir, las variables $x, y, z$ entran de la misma forma en la ecuación), los cálculos para las otras coordenadas son análogos:

$$ \frac{\partial r}{\partial y} = \frac{y}{(x^2 + y^2 + z^2)^{1/2}} = \frac{y}{r} $$
$$ \frac{\partial r}{\partial z} = \frac{z}{(x^2 + y^2 + z^2)^{1/2}} = \frac{z}{r} $$

\subsection{Parte (b): Gradiente y verificación de magnitud}

\subsubsection{Demostración de $\nabla r = \frac{\vec{r}}{r}$}
El operador gradiente $\nabla$ en un sistema cartesiano se define como el vector cuyas componentes son las derivadas parciales de la función escalar sobre la cual actúa:
$$ \nabla r = \left( \frac{\partial r}{\partial x}, \frac{\partial r}{\partial y}, \frac{\partial r}{\partial z} \right) $$

Sustituyendo los resultados de la parte (a):

$$ \nabla r = \left( \frac{x}{r}, \frac{y}{r}, \frac{z}{r} \right) $$

Podemos extraer el factor escalar común $\frac{1}{r}$:
$$ \nabla r = \frac{1}{r} (x, y, z) $$

Dado que el vector $(x, y, z)$ es la definición de $\vec{r}$ en componentes cartesianas, llegamos a la conclusión requerida:

$$ \nabla r = \frac{\vec{r}}{r} = \hat{r} $$

Donde $\hat{r}$ es el vector unitario en la dirección radial.

\subsubsection{Verificación de que $|\nabla r| = 1$}
Para verificar la magnitud del vector resultante, calculamos su norma euclidiana:
$$ |\nabla r| = \sqrt{ \left( \frac{x}{r} \right)^2 + \left( \frac{y}{r} \right)^2 + \left( \frac{z}{r} \right)^2 } $$

Elevando al cuadrado los términos:

$$ |\nabla r| = \sqrt{ \frac{x^2}{r^2} + \frac{y^2}{r^2} + \frac{z^2}{r^2} } $$
$$|\nabla r| = \sqrt{ \frac{x^2 + y^2 + z^2}{r^2} } $$

Utilizando la definición original de $r^2 = x^2 + y^2 + z^2$:

$$ |\nabla r| = \sqrt{ \frac{r^2}{r^2} } $$
$$ |\nabla r| = \sqrt{1} = 1 $$
Este resultado es válido siempre que $r \neq 0$ para evitar la indeterminación en el origen.

\section*{Relevancia en física}

Este problema es fundamental en la física por varias razones clave:

\begin{enumerate}
    \item \textbf{Fuerzas Centrales:} En Mecánica Celeste y Electrodinámica, muchas fuerzas dependen solo de la distancia $r$. Este cálculo permite demostrar que la fuerza es proporcional al gradiente del potencial radial, $\vec{F} = -\nabla V(r) = -\hat{r} \frac{dV}{dr}$, lo que implica que las fuerzas gravitacionales y electrostáticas siempre apuntan a lo largo de la línea que une a los cuerpos.
    \item \textbf{Propiedades del Espacio Isotrópico:} El hecho de que $|\nabla r| = 1$ refleja la homogeneidad e isotropía del espacio euclidiano, donde la tasa de cambio de la distancia respecto a un desplazamiento en su propia dirección es la unidad.
    \item \textbf{Cálculo de Potenciales:} Es la base para derivar campos de vectores a partir de potenciales escalares (campos conservativos), lo cual simplifica enormemente la resolución de problemas de dinámica $N$-cuerpos.
\end{enumerate}

\section{Problema \# 2}

En coordenadas cartesianas, sea la función escalar $f(\vec{r}) = \frac{1}{r}$, donde $r = |\vec{r}| = (x^2 + y^2 + z^2)^{1/2}$ es la magnitud del vector de posición $\vec{r} = (x, y, z)$.

\begin{enumerate}[a)]
    \item Calcule la derivada parcial $\frac{\partial}{\partial x} \left( \frac{1}{r} \right)$. Sugerencia: use la definición de $r$ en términos de sus componentes cartesianas.
    \item Construya el vector gradiente $\nabla \left( \frac{1}{r} \right)$ y demuestre que, para todo $r \neq 0$, se cumple que $\nabla \left( \frac{1}{r} \right) = -\frac{\vec{r}}{r^3}$.
\end{enumerate}

\section*{Solución Problema \# 2}

Este problema aborda el cálculo del \textbf{gradiente de un potencial escalar inversamente proporcional a la distancia}, una operación fundamental en la física clásica. El objetivo es demostrar matemáticamente cómo la variación de la función $1/r$ en el espacio tridimensional genera un campo vectorial que apunta hacia el origen y cuya magnitud decae con el cuadrado de la distancia ($1/r^2$).

\subsection{Parte (a): Cálculo de la derivada parcial}

Para calcular $\frac{\partial}{\partial x} \left( \frac{1}{r} \right)$, tratamos a $r$ como una función compuesta de $x, y, z$ y aplicamos la \textbf{regla de la cadena}.

\subsubsection{Derivación paso a paso}

Definimos la función como $f(r) = r^{-1}$. La derivada parcial respecto a $x$ es:

$$ \frac{\partial f}{\partial x} = \frac{df}{dr} \frac{\partial r}{\partial x} $$

Primero, calculamos la derivada de $f$ respecto a $r$:
$$ \frac{df}{dr} = \frac{d}{dr} (r^{-1}) = -r^{-2} = -\frac{1}{r^2} $$

Segundo, utilizamos el resultado del problema anterior (o lo derivamos nuevamente) para la variación de $r$ respecto a $x$:
$$ \frac{\partial r}{\partial x} = \frac{\partial}{\partial x} (x^2 + y^2 + z^2)^{1/2} = \frac{1}{2}(x^2 + y^2 + z^2)^{-1/2} \cdot (2x) = \frac{x}{r} $$

Finalmente, combinamos ambos resultados:
$$ \frac{\partial}{\partial x} \left( \frac{1}{r} \right) = \left( -\frac{1}{r^2} \right) \left( \frac{x}{r} \right) = -\frac{x}{r^3} $$

\subsection{Parte (b): Construcción del gradiente y demostración}

El gradiente $\nabla$ en coordenadas cartesianas de una función escalar $\phi$ se define como el vector:
$ \nabla \phi = \left( \frac{\partial \phi}{\partial x}, \frac{\partial \phi}{\partial y}, \frac{\partial \phi}{\partial z} \right) $

\subsubsection{Cálculo de las componentes restantes}

Dada la simetría de la función $1/r$, las derivadas respecto a $y$ y $z$ siguen exactamente la misma forma que la obtenida para $x$:
$$ \frac{\partial}{\partial y} \left( \frac{1}{r} \right) = -\frac{y}{r^3} $$
$$ \frac{\partial}{\partial z} \left( \frac{1}{r} \right) = -\frac{z}{r^3} $$

\subsubsection{Ensamblaje del vector gradiente}

Sustituyendo las tres componentes en la definición del gradiente:
$$ \nabla \left( \frac{1}{r} \right) = \left( -\frac{x}{r^3}, -\frac{y}{r^3}, -\frac{z}{r^3} \right) $$

Podemos factorizar el término escalar común $-\frac{1}{r^3}$:
$$ \nabla \left( \frac{1}{r} \right) = -\frac{1}{r^3} (x, y, z) $$

Como el vector $(x, y, z)$ es por definición el vector de posición $\vec{r}$, llegamos al resultado final:
$$ \nabla \left( \frac{1}{r} \right) = -\frac{\vec{r}}{r^3} $$

Para expresar esto en términos del vector unitario $\hat{r} = \frac{\vec{r}}{r}$, también podemos escribir:
$$ \nabla \left( \frac{1}{r} \right) = -\frac{1}{r^2} \hat{r} $$

\section*{Relevancia en física}

Este resultado es de vital importancia en la \textbf{Mecánica Celeste} y la física general por las siguientes razones:

\begin{enumerate}
    \item \textbf{Potencial Gravitacional:} La ley de gravitación de Newton establece que la fuerza es central y conservativa. El potencial gravitacional de una masa puntual $M$ es $\Phi = -G \frac{M}{r}$. La aceleración de la gravedad se obtiene como $\vec{g} = -\nabla \Phi$. Al aplicar el resultado de este problema, vemos que $\vec{g} = -GM \nabla(1/r) = -G \frac{M}{r^3} \vec{r}$, que es la forma vectorial de la \textbf{ley del inverso del cuadrado}.
    \item \textbf{Ecuación de Poisson y Laplace:} Este cálculo demuestra que $1/r$ es una solución a la ecuación de Laplace ($\nabla^2 \phi = 0$) en todo punto excepto en el origen ($r=0$). De hecho, en el origen, este gradiente genera una singularidad relacionada con la función delta de Dirac, esencial para entender fuentes puntuales de masa o carga.
    \item \textbf{Campos Conservativos:} El hecho de que una fuerza pueda escribirse como el gradiente de un potencial escalar ($1/r$) garantiza la conservación de la energía mecánica en el sistema de dos cuerpos.
\end{enumerate}

\section{Problema \# 3}

(En coordenadas esféricas) Use el operador gradiente en coordenadas esféricas para una función escalar $\phi(r, \theta, \varphi)$, definido como:
$$ \nabla \phi = \frac{\partial \phi}{\partial r} \hat{e}_r + \frac{1}{r} \frac{\partial \phi}{\partial \theta} \hat{e}_\theta + \frac{1}{r \sin \theta} \frac{\partial \phi}{\partial \varphi} \hat{e}_\varphi $$

\begin{enumerate}[a)]
    \item Para la función $\phi(r, \theta, \varphi) = \frac{1}{r}$, calcule $\nabla \phi$.
    \item Compare con el resultado obtenido en el problema 2 (en coordenadas cartesianas) y explique por qué una función puramente radial simplifica el cálculo en problemas de fuerzas centrales.
\end{enumerate}

\section*{Solución Problema \# 3}

Este problema explora el cálculo del gradiente utilizando el sistema de \textbf{coordenadas esféricas}, el cual es el sistema de referencia natural para problemas con simetría central. El objetivo es verificar que, para un potencial radial como $1/r$, el uso de coordenadas esféricas reduce significativamente la complejidad algebraica en comparación con las coordenadas cartesianas, facilitando la descripción de campos de fuerza como el gravitatorio.

\subsection{Parte (a): Cálculo del gradiente en esféricas}

Dada la función escalar $\phi(r, \theta, \varphi) = \frac{1}{r}$, observamos que esta función depende exclusivamente de la coordenada radial $r$ y es independiente de los ángulos $\theta$ (colatitud) y $\varphi$ (azimut).

\subsubsection{Derivadas parciales}

Procedemos a calcular cada componente del gradiente según la definición proporcionada:

1. \textbf{Componente radial:}

$$ \frac{\partial \phi}{\partial r} = \frac{\partial}{\partial r} \left( \frac{1}{r} \right) = -\frac{1}{r^2} $$

2. \textbf{Componente polar ($\theta$):}
Como $\phi$ no contiene la variable $\theta$:
$$ \frac{1}{r} \frac{\partial \phi}{\partial \theta} = \frac{1}{r} (0) = 0 $$

3. \textbf{Componente azimutal ($\varphi$):}
Como $\phi$ no contiene la variable $\varphi$:
$$ \frac{1}{r \sin \theta} \frac{\partial \phi}{\partial \varphi} = \frac{1}{r \sin \theta} (0) = 0 $$

\subsubsection{Resultado final}

Sustituyendo estos valores en la expresión del gradiente:
$$ \nabla \left( \frac{1}{r} \right) = -\frac{1}{r^2} \hat{e}_r + 0 \hat{e}_\theta + 0 \hat{e}_\varphi $$
$$ \nabla \left( \frac{1}{r} \right) = -\frac{1}{r^2} \hat{e}_r $$

\subsection{Parte (b): Comparación y simplificación en problemas centrales}

\subsubsection{Comparación con coordenadas cartesianas}
En el problema 2, se demostró en coordenadas cartesianas que $\nabla (1/r) = -\frac{\vec{r}}{r^3}$. 
Para comparar, recordamos que el vector de posición se define como $\vec{r} = r \hat{e}_r$. Al sustituir esto en el resultado cartesiano obtenemos:
$$ -\frac{\vec{r}}{r^3} = -\frac{r \hat{e}_r}{r^3} = -\frac{1}{r^2} \hat{e}_r $$
Ambos resultados son idénticos, lo que confirma la \textbf{invarianza del operador gradiente} bajo transformaciones de coordenadas.

\subsubsection{¿Por qué simplifica el cálculo?}
El uso de una función puramente radial simplifica el cálculo en problemas de fuerzas centrales por las siguientes razones:
\begin{itemize}
    \item \textbf{Anulación de términos:} Al depender solo de $r$, las derivadas respecto a las coordenadas angulares ($\theta, \varphi$) desaparecen automáticamente, reduciendo el gradiente a un solo término.
    \item \textbf{Alineación natural:} El resultado apunta directamente en la dirección del vector unitario radial $\hat{e}_r$, lo cual coincide con la línea de acción de la fuerza gravitacional.
    \item \textbf{Eficiencia algebraica:} Evita la manipulación de raíces cuadradas y potencias de componentes $(x, y, z)$ que suelen aparecer en cartesianas, permitiendo una transición directa del potencial escalar a la aceleración vectorial.
\end{itemize}

\section*{Relevancia en física}

Este ejercicio tiene una relevancia fundamental en la {Mecánica Celeste}:

\begin{enumerate}
    \item \textbf{Leyes de Kepler:} La mayoría de los cuerpos celestes se modelan como masas puntuales donde el potencial gravitatorio depende únicamente de la distancia $r$. Las coordenadas esféricas permiten derivar las ecuaciones de movimiento de forma mucho más directa, facilitando la deducción de las órbitas cónicas.
    \item \textbf{Campos de Fuerza Central:} Una fuerza es 'central' si su magnitud depende solo de $r$ y su dirección es radial. Este problema demuestra matemáticamente que cualquier potencial escalar de la forma $V(r)$ genera una fuerza central $\vec{F} = -\nabla V = -\frac{dV}{dr} \hat{e}_r$, lo que garantiza la conservación del momento angular.
    \item \textbf{Estabilidad Numérica:} En astrodinámica computacional, entender estas propiedades permite elegir sistemas de coordenadas que minimizan los errores de redondeo y simplifican la integración de trayectorias complejas.
\end{enumerate}

\section{Problema \# 4}

Sean $\vec{r}(t)$ y $\vec{v}(t)$ los vectores de posición y velocidad de una partícula, donde $\vec{v}(t) = \dot{\vec{r}}(t)$. Definimos la magnitud del vector de posición como $r = |\vec{r}|$ y el vector unitario en la dirección radial como $\hat{r} = \frac{\vec{r}}{r}$.

\begin{enumerate}[a)]
    \item Demuestre que la derivada temporal de la magnitud del vector de posición (velocidad radial) está dada por $\dot{r} = \hat{r} \cdot \vec{v}$.
    \item A partir de lo anterior, demuestre que la derivada temporal del vector unitario radial es $\frac{d\hat{r}}{dt} = \frac{\vec{v} - \dot{r}\hat{r}}{r}$.
    \item Justifique algebraicamente por qué el producto escalar $\hat{r} \cdot \frac{d\hat{r}}{dt}$ es siempre igual a cero.
\end{enumerate}

\section*{Solución Problema \# 4}

Este problema es fundamental para la cinemática en coordenadas no cartesianas (como polares o esféricas). Se centra en entender cómo cambian tanto la **longitud** como la **dirección** del vector de posición a medida que una partícula se mueve. El objetivo es descomponer el vector velocidad en una componente que cambia la distancia al origen y otra que cambia la orientación del vector unitario radial.

\subsection{Parte (a): Derivada de la magnitud $r$}

Para encontrar la tasa de cambio de la distancia $r$, partimos de la definición del cuadrado de la magnitud usando el producto escalar:
$$ r^2 = \vec{r} \cdot \vec{r} $$

Derivamos ambos lados de la ecuación con respecto al tiempo $t$ aplicando la regla del producto:
$$ \frac{d}{dt}(r^2) = \frac{d}{dt}(\vec{r} \cdot \vec{r}) $$
$$ 2r \frac{dr}{dt} = \dot{\vec{r}} \cdot \vec{r} + \vec{r} \cdot \dot{\vec{r}} $$

Dado que el producto escalar es conmutativo y definiendo $\dot{r} = \frac{dr}{dt}$ y $\vec{v} = \dot{\vec{r}}$:
$$ 2r\dot{r} = 2\vec{r} \cdot \vec{v} $$

Dividiendo por $2r$ (asumiendo $r \neq 0$):
$$ \dot{r} = \frac{\vec{r} \cdot \vec{v}}{r} $$

Sustituyendo la definición de vector unitario $\hat{r} = \frac{\vec{r}}{r}$, llegamos al resultado:
$$ \dot{r} = \hat{r} \cdot \vec{v} $$

Este resultado indica que la velocidad radial es la proyección del vector velocidad total sobre la dirección radial.

\subsection{Parte (b): Derivada del vector unitario $\hat{r}$}

Para calcular $\frac{d\hat{r}}{dt}$, aplicamos la regla del cociente a la definición $\hat{r} = \frac{\vec{r}}{r}$, tratando a $\vec{r}$ como el numerador y a $r$ como el denominador:
$$ \frac{d\hat{r}}{dt} = \frac{d}{dt} \left( \frac{\vec{r}}{r} \right) = \frac{\dot{\vec{r}}r - \vec{r}\dot{r}}{r^2} $$

Sustituimos $\dot{\vec{r}}$ por $\vec{v}$:
$$ \frac{d\hat{r}}{dt} = \frac{\vec{v}r - \vec{r}\dot{r}}{r^2} $$

Podemos separar la fracción para simplificar los términos:
$$ \frac{d\hat{r}}{dt} = \frac{\vec{v}r}{r^2} - \frac{\vec{r}\dot{r}}{r^2} = \frac{\vec{v}}{r} - \frac{\dot{r}}{r} \left( \frac{\vec{r}}{r} \right) $$

Reconocemos nuevamente que $\frac{\vec{r}}{r} = \hat{r}$, por lo que:
$$ \frac{d\hat{r}}{dt} = \frac{\vec{v} - \dot{r}\hat{r}}{r} $$

Este vector representa la tasa de cambio en la dirección del vector de posición. Es importante notar que el numerador $\vec{v} - \dot{r}\hat{r}$ es la componente de la velocidad perpendicular a la dirección radial.

\subsection{Parte (c): Justificación de $\hat{r} \cdot \dot{\hat{r}} = 0$}

Existen dos formas de justificar este resultado: una geométrica y otra algebraica.

\subsubsection{Justificación algebraica (Propiedad de la norma constante)}
Un vector unitario tiene, por definición, una magnitud constante igual a 1. Por lo tanto:
$$ \hat{r} \cdot \hat{r} = |\hat{r}|^2 = 1 $$

Si derivamos esta constante respecto al tiempo:
$$ \frac{d}{dt}(\hat{r} \cdot \hat{r}) = \frac{d}{dt}(1) $$
$$ \dot{\hat{r}} \cdot \hat{r} + \hat{r} \cdot \dot{\hat{r}} = 0 $$
$$ 2\hat{r} \cdot \dot{\hat{r}} = 0 \implies \hat{r} \cdot \dot{\hat{r}} = 0 $$

\subsubsection{Interpretación física}
Este resultado demuestra que el cambio de un vector unitario siempre debe ser perpendicular al vector mismo. Si tuviera una componente en la dirección radial, cambiaría la longitud del vector, lo que contradice su definición como vector unitario.

\section*{Relevancia en física}

Este problema es la base técnica para construir la cinemática en sistemas de referencia rotantes y es vital en la **mecánica celeste**:

\begin{enumerate}
    \item \textbf{Momento Angular:} La relación $\dot{\hat{r}}$ permite demostrar que en un sistema de dos cuerpos, la velocidad transversal es responsable del barrido de áreas (Segunda Ley de Kepler).
    \item \textbf{Aceleraciones Centrales:} Al derivar nuevamente estos resultados, aparecen los términos de aceleración centrípeta y de Coriolis, esenciales para el Problema de los $N$-Cuerpos.
    \item \textbf{Elementos Orbitales:} La determinación de la anomalía verdadera $f$ y su evolución temporal depende directamente de entender cómo la velocidad total se reparte entre el cambio en la distancia ($r$) y el cambio en la orientación ($\hat{r}$).
\end{enumerate}


\section{Problema \# 5}

Sea $U(\vec{r})$ un potencial escalar que depende únicamente de la posición $\vec{r}$, y sea $\vec{r}(t)$ una trayectoria definida en el espacio tridimensional en función del tiempo $t$.

\begin{enumerate}[a)]
    \item Escriba la derivada total $\frac{dU}{dt}$ utilizando la regla de la cadena multivariada en términos de las derivadas parciales $\frac{\partial U}{\partial x}$, $\frac{\partial U}{\partial y}$, $\frac{\partial U}{\partial z}$ y las componentes de la velocidad $\dot{x}, \dot{y}, \dot{z}$.
    \item Reescriba el resultado utilizando notación vectorial y demuestre que $\frac{dU}{dt} = \nabla U(\vec{r}(t)) \cdot \dot{\vec{r}}(t)$.
\end{enumerate}

\section*{Solución Problema \# 5}

Este problema aborda el concepto de la **tasa de cambio temporal** de un campo escalar (en este caso, un potencial gravitatorio) cuando este es "visto" por una partícula que se desplaza a través del espacio siguiendo una trayectoria determinada. El objetivo es comprender cómo la variación espacial del potencial (gradiente) se proyecta sobre el movimiento de la partícula (velocidad) para producir un cambio neto en el valor del potencial a lo largo del tiempo.

\subsection{Parte (a): Aplicación de la regla de la cadena multivariada}

Dado que el potencial $U$ depende de la posición $\vec{r} = (x, y, z)$, y a su vez cada una de estas coordenadas es una función del tiempo $t$ debido al movimiento de la partícula ($x(t), y(t), z(t)$), podemos expresar a $U$ como una función compuesta del tiempo: $U(x(t), y(t), z(t))$.

Para encontrar la derivada total de $U$ respecto a $t$, aplicamos la **regla de la cadena para funciones de varias variables**. Esta establece que el cambio total es la suma de los cambios producidos a través de cada una de sus variables independientes:

$$ \frac{dU}{dt} = \frac{\partial U}{\partial x} \frac{dx}{dt} + \frac{\partial U}{\partial y} \frac{dy}{dt} + \frac{\partial U}{\partial z} \frac{dz}{dt} $$

Si adoptamos la notación de punto para las derivadas temporales, donde $\dot{x} = \frac{dx}{dt}$, $\dot{y} = \frac{dy}{dt}$ y $\dot{z} = \frac{dz}{dt}$, la expresión anterior se escribe de forma más sucinta como:

$$ \frac{dU}{dt} = \frac{\partial U}{\partial x} \dot{x} + \frac{\partial U}{\partial y} \dot{y} + \frac{\partial U}{\partial z} \dot{z} $$

Es importante notar que, dado que el enunciado especifica que el potencial depende **solo de la posición**, no existe un término de variación temporal explícita ($\frac{\partial U}{\partial t} = 0$), por lo que el cambio se debe exclusivamente al desplazamiento de la partícula en el campo espacial.

\subsection{Parte (b): Notación vectorial y demostración}

Para simplificar la expresión obtenida, podemos identificar las entidades vectoriales subyacentes en el cálculo anterior.

\subsubsection{Identificación de los vectores de estado}

1. \textbf{El Gradiente del Potencial ($\nabla U$):} Se define como el vector cuyas componentes son las derivadas parciales de la función escalar respecto a las coordenadas espaciales:
$$ \nabla U = \left( \frac{\partial U}{\partial x}, \frac{\partial U}{\partial y}, \frac{\partial U}{\partial z} \right) $$

2. \textbf{El Vector Velocidad ($\dot{\vec{r}}$):} Es el vector que representa la tasa de cambio de la posición de la partícula en el tiempo:
$$ \dot{\vec{r}} = (\dot{x}, \dot{y}, \dot{z}) $$

\subsubsection{Demostración mediante el producto escalar}

Recordando la definición del **producto escalar** entre dos vectores $\vec{A} \cdot \vec{B} = \sum_{i} A_i B_i$, observamos que la estructura de la derivada total obtenida en el inciso (a) coincide exactamente con esta operación matemática:

$$ \frac{dU}{dt} = \left( \frac{\partial U}{\partial x}, \frac{\partial U}{\partial y}, \frac{\partial U}{\partial z} \right) \cdot (\dot{x}, \dot{y}, \dot{z}) $$

Por lo tanto, llegamos a la identidad vectorial requerida:

$$ \frac{dU}{dt} = \nabla U(\vec{r}(t)) \cdot \dot{\vec{r}}(t) $$

Esta expresión compacta nos dice que la tasa de cambio del potencial es igual a la proyección del gradiente del campo sobre el vector velocidad de la partícula.

\section*{Relevancia en física}

Este resultado es una piedra angular tanto en la física general como en la **mecánica celeste** por las siguientes razones fundamentales:

\begin{enumerate}
    \item \textbf{Conservación de la Energía:} Este cálculo es el paso matemático necesario para demostrar que la energía mecánica se conserva. Al multiplicar la ecuación de movimiento ($m\vec{a} = \vec{F}$) por la velocidad, el término de la fuerza conservativa ($\vec{F} \cdot \vec{v} = -\nabla U \cdot \vec{v}$) se convierte, gracias a este problema, en $-dU/dt$, permitiendo la integración directa de la energía.
    \item \textbf{Trabajo y Potencia:} Define la potencia mecánica de una fuerza conservativa. La potencia es $P = \vec{F} \cdot \vec{v}$; si la fuerza es central y conservativa, la potencia es simplemente el negativo de la derivada temporal del potencial.
    \item \textbf{Dinámica en Sistemas de N-Cuerpos:} Permite describir la evolución de sistemas complejos (como el Sistema Solar) mediante un solo funcional escalar, facilitando el uso de formalismos avanzados como el Lagrangiano o el Hamiltoniano.
\end{enumerate}


\section{Problema \# 6}

En un sistema de $N$ cuerpos, defina la distancia relativa entre las partículas $j$ e $i$ como $r_{ji} = |\vec{r}_j - \vec{r}_i|$.

\begin{enumerate}[a)]
    \item Calcule el gradiente de la función escalar $\frac{1}{r_{ji}}$ con respecto a las coordenadas de la partícula $j$, denotado como $\nabla_j \left( \frac{1}{r_{ji}} \right)$.
    \item Calcule el gradiente de la misma función con respecto a las coordenadas de la partícula $i$, denotado como $\nabla_i \left( \frac{1}{r_{ji}} \right)$.
    \item Establezca la relación exacta entre ambos resultados y exprésela de forma compacta.
\end{enumerate}

\section*{Solución Problema \# 6}

Este problema extiende el análisis del gradiente de un potencial central al contexto de un sistema de múltiples partículas puntuales. El objetivo es comprender cómo varía la energía potencial de interacción mutua (proporcional a $1/r_{ji}$) cuando se desplaza individualmente a cada uno de los cuerpos involucrados en el par. Este análisis es fundamental para derivar las fuerzas internas en el problema de los $N$-cuerpos y verificar la consistencia con las leyes de Newton.

\subsection{Parte (a): Gradiente respecto a la partícula $j$}

Definimos el vector de posición relativa que apunta de $i$ hacia $j$ como $\vec{r}_{ji} = \vec{r}_j - \vec{r}_i$. Al calcular el gradiente $\nabla_j$, tratamos a $\vec{r}_i$ como un vector constante.

\subsubsection{Cálculo paso a paso}

La variación de la magnitud $r_{ji}$ respecto a un cambio en la posición de la partícula $j$ es idéntica al gradiente radial simple estudiado en problemas anteriores:
$$ \nabla_j r_{ji} = \frac{\vec{r}_{ji}}{r_{ji}} = \hat{r}_{ji} $$

Aplicando la regla de la cadena para la función recíproca:
$$ \nabla_j \left( \frac{1}{r_{ji}} \right) = \frac{d}{dr_{ji}} \left( r_{ji}^{-1} \right) \nabla_j r_{ji} $$
$$ \nabla_j \left( \frac{1}{r_{ji}} \right) = \left( -\frac{1}{r_{ji}^2} \right) \left( \frac{\vec{r}_{ji}}{r_{ji}} \right) $$
$$ \nabla_j \left( \frac{1}{r_{ji}} \right) = -\frac{\vec{r}_{ji}}{r_{ji}^3} $$

Este resultado indica que el gradiente respecto a la partícula que "siente" el potencial apunta hacia la otra partícula $,$.

\subsection{Parte (b): Gradiente respecto a la partícula $i$}

Ahora calculamos el gradiente $\nabla_i$, tratando a $\vec{r}_j$ como constante. Notamos que en la definición $r_{ji} = \sqrt{(x_j-x_i)^2 + (y_j-y_i)^2 + (z_j-z_i)^2}$, las coordenadas de la partícula $i$ entran con signo negativo.

\subsubsection{Cálculo de la derivada interna}

Calculamos la derivada parcial de $r_{ji}$ respecto a $x_i$:
$$ \frac{\partial r_{ji}}{\partial x_i} = \frac{1}{2 r_{ji}} \cdot 2(x_j - x_i) \cdot (-1) = -\frac{x_j - x_i}{r_{ji}} = \frac{x_i - x_j}{r_{ji}} $$

Al generalizar para las tres componentes:
$$ \nabla_i r_{ji} = \frac{\vec{r}_i - \vec{r}_j}{r_{ji}} = -\frac{\vec{r}_{ji}}{r_{ji}} $$

Finalmente, el gradiente de la función recíproca es:
$$ \nabla_i \left( \frac{1}{r_{ji}} \right) = \left( -\frac{1}{r_{ji}^2} \right) \nabla_i r_{ji} = \left( -\frac{1}{r_{ji}^2} \right) \left( -\frac{\vec{r}_{ji}}{r_{ji}} \right) $$
$$ \nabla_i \left( \frac{1}{r_{ji}} \right) = \frac{\vec{r}_{ji}}{r_{ji}^3} $$

\subsection{Parte (c): Relación entre gradientes}

Comparando los resultados obtenidos en (a) y (b), observamos que los vectores resultantes tienen la misma magnitud y dirección, pero sentidos opuestos:
$$ \nabla_j \left( \frac{1}{r_{ji}} \right) = -\frac{\vec{r}_{ji}}{r_{ji}^3} \quad , \quad \nabla_i \left( \frac{1}{r_{ji}} \right) = \frac{\vec{r}_{ji}}{r_{ji}^3} $$

La relación compacta es:
$$ \nabla_j \left( \frac{1}{r_{ji}} \right) = -\nabla_i \left( \frac{1}{r_{ji}} \right) $$

Esta simetría es una propiedad intrínseca de las funciones que dependen únicamente de la distancia relativa entre dos puntos $,$.

\section*{Relevancia en física}

Este resultado es de suma importancia en la **Mecánica Celeste** y la dinámica de sistemas por las siguientes razones:

\begin{enumerate}
    \item \textbf{Tercera Ley de Newton:} La fuerza ejercida por la partícula $j$ sobre la partícula $i$ es $\vec{F}_{ij} = -G m_i m_j \nabla_i (1/r_{ji})$. Debido a la relación demostrada, se cumple automáticamente que $\vec{F}_{ij} = -\vec{F}_{ji}$, satisfaciendo el postulado de acción y reacción.
    \item \textbf{Conservación del Momentum Lineal:} Esta simetría en las derivadas del potencial garantiza que la suma de las fuerzas internas en un sistema de $N$ cuerpos sea nula, lo que implica que el centro de masa se mueve de forma inercial.
    \item \textbf{Formulación del Hamiltoniano:} En la mecánica analítica, permite simplificar el cálculo de las ecuaciones de movimiento al reducir el número de operaciones necesarias para evaluar las interacciones gravitatorias mutuas.
\end{enumerate}




\bibliographystyle{plainurl}
\bibliography{bibliografia} 
\end{document}
